\documentclass[12pt,a4paper]{article}
\usepackage[utf8]{inputenc}
\usepackage[T1]{fontenc}
\usepackage[turkish]{babel}
\usepackage{geometry}
\usepackage{booktabs}
\usepackage{longtable}
\usepackage{array}
\usepackage{hyperref}
\usepackage{xcolor}
\usepackage{titlesec}
\usepackage{enumitem}
\usepackage{fancyhdr}
\usepackage{tabularx}
\usepackage{lmodern}

\geometry{margin=2.5cm}

\definecolor{primary}{RGB}{88, 56, 156}
\definecolor{accent}{RGB}{147, 112, 219}
\definecolor{lightgray}{RGB}{245, 245, 245}

\hypersetup{
  colorlinks=true,
  linkcolor=primary,
  urlcolor=primary
}

\titleformat{\section}{\large\bfseries\color{primary}}{}{0em}{}[\titlerule]
\titleformat{\subsection}{\normalsize\bfseries\color{accent}}{}{0em}{}

\pagestyle{fancy}
\fancyhf{}
\rhead{\textcolor{gray}{\small CanYoldaşı — Geliştirme Durum Raporu}}
\lhead{\textcolor{gray}{\small 16 Temmuz 2026}}
\rfoot{\textcolor{gray}{\small \thepage}}

\begin{document}

% ── Başlık ───────────────────────────────────────────────────────────────────
\begin{center}
  {\LARGE\bfseries\color{primary} CanYoldaşı}\\[0.4em]
  {\large Geliştirme Durum Raporu}\\[0.3em]
  {\normalsize\color{gray} 16 Temmuz 2026}
\end{center}

\vspace{1em}
\noindent
CanYoldaşı, Türkiye'deki sokak hayvanlarını raporlamak, takip etmek ve sahiplendirme ilanları
yönetmek için geliştirilen topluluk odaklı bir mobil uygulamadır.
\vspace{1em}

% ── Stack ────────────────────────────────────────────────────────────────────
\section{Teknik Altyapı}

\begin{tabularx}{\textwidth}{lX}
  \toprule
  \textbf{Katman} & \textbf{Teknoloji} \\
  \midrule
  Mobil & Expo SDK 54, expo-router, React Native 0.81 \\
  Backend & Express 5 + PostgreSQL + Drizzle ORM \\
  Kimlik Doğrulama & JWT (email/şifre, AsyncStorage) \\
  Harita & react-native-maps@1.18.0 (sabitlenmiş) \\
  IAP Ödemesi & RevenueCat (iOS + Android) \\
  Web Ödemesi & Stripe \\
  E-posta & Resend (şifre sıfırlama) \\
  Monorepo & pnpm workspaces, TypeScript 5.9 \\
  \bottomrule
\end{tabularx}

% ── Tamamlanan Adımlar ───────────────────────────────────────────────────────
\section{Tamamlanan Adımlar}

% ── Adım 10–11 ──
\subsection{Adım 10–11 — RevenueCat IAP Entegrasyonu}

\begin{itemize}[leftmargin=1.5em]
  \item \texttt{react-native-purchases} SDK entegrasyonu (iOS + Android)
  \item Ürün tanımları: \texttt{canyoldasi\_boost\_1\_day}, \texttt{canyoldasi\_boost\_3\_days},
        \texttt{canyoldasi\_boost\_7\_days}
  \item IAP satın alma akışı: platform ödeme sayfası $\rightarrow$ RevenueCat $\rightarrow$ sunucu doğrulama
\end{itemize}

% ── Adım 12 ──
\subsection{Adım 12 — Verify-Purchase Atomik Pipeline}

Endpoint: \texttt{POST /api/promotions/verify-purchase}

\begin{tabularx}{\textwidth}{lX}
  \toprule
  \textbf{Güvenlik Katmanı} & \textbf{Detay} \\
  \midrule
  JWT zorunlu & 401 — anonim istekler reddedilir \\
  İlan sahipliği & \texttt{listing.userId === auth userId} — 403 \\
  \texttt{durationDays} & Sunucu tarafı \texttt{PACKAGE\_MAP}'ten; istemciden alınmaz \\
  Stacking & Aktif promosyon varsa üstüne eklenir \\
  Idempotency & \texttt{transaction\_identifier} UNIQUE index — çift gönderim güvenli \\
  \bottomrule
\end{tabularx}

% ── Adım 13 ──
\subsection{Adım 13 — Promosyon UI ve Paylaşılan Yardımcılar}

\begin{itemize}[leftmargin=1.5em]
  \item \texttt{promotionHelpers.ts}: \texttt{isListingPromoted()}, \texttt{formatRemainingTime()},
        \texttt{useNow()} (dakikalık tick)
  \item Mor/leylak badge: \textit{⭐ Öne Çıkan · 2g 14s}
  \item ``Öne Çıkarmayı Uzat'' stacking CTA butonu
  \item Öne çıkan ilanlar liste başında sıralanır
  \item iOS 26+ Liquid Glass sekme çubuğu; eski sürümlerde BlurView fallback
\end{itemize}

% ── Adım 14 ──
\subsection{Adım 14 — RevenueCat Webhook + İade/İptal Yaşam Döngüsü}

Yeni endpoint: \texttt{POST /api/webhooks/revenuecat}

\vspace{0.5em}
\noindent\textbf{Güvenlik Katmanları (sırayla):}

\begin{enumerate}[leftmargin=1.5em]
  \item \texttt{Authorization} header $\leftrightarrow$ \texttt{REVENUECAT\_WEBHOOK\_AUTH\_TOKEN}
        ortam değişkeni — \texttt{crypto.timingSafeEqual} ile sabit zamanlı karşılaştırma;
        ortam değişkeni yoksa fail-closed (tüm istekler reddedilir)
  \item \texttt{X-RevenueCat-Webhook-Signature: t=\textless ts\textgreater,v1=\textless hex\textgreater}
        HMAC-SHA256 $\leftrightarrow$ \texttt{REVENUECAT\_WEBHOOK\_SIGNING\_SECRET}
        — opsiyonel; set edilmişse zorunludur; 5 dakikadan eski imzalar reddedilir
  \item Zod şema doğrulama
  \item \texttt{revenuecat\_event\_id} UNIQUE kontrolü — çift teslimat idempotent 200 döner
\end{enumerate}

\vspace{0.5em}
\noindent\textbf{Event Yönlendirme Tablosu:}

\begin{tabularx}{\textwidth}{lX}
  \toprule
  \textbf{Event Türü} & \textbf{Davranış} \\
  \midrule
  \texttt{NON\_RENEWING\_PURCHASE} &
    Webhook reconciliation yolu — \texttt{verified} durumundaki satın alımı
    \texttt{processed}'e geçirir ve ilanı aktive eder \\
  \texttt{CANCELLATION} &
    \texttt{processed} $\rightarrow$ \texttt{refunded}; promosyon ledger'ı
    yeniden hesaplanır \\
  \texttt{REFUND\_REVERSED} &
    \texttt{refunded} $\rightarrow$ \texttt{processed} restore edilir;
    promosyon ledger'ı yeniden hesaplanır \\
  \texttt{TEST} &
    Kaydedilir, \texttt{skipped} olarak işaretlenir; ilan mutasyonu yapılmaz \\
  bilinmeyen & Güvenli \texttt{skipped} — uygulama crash olmaz \\
  \bottomrule
\end{tabularx}

\vspace{0.5em}
\noindent\textbf{Promosyon Ledger Algoritması (\texttt{recalculateListingPromotion}):}

\vspace{0.3em}
\noindent
İade sonrasında kalan tüm \texttt{purchaseStatus = "processed"} satın alımlar üzerinden
$\max(\texttt{promotionExpiresAt})$ hesaplanır ve \texttt{adoption\_listings.promoted\_until}
güncellenir. Stacking senaryolarında geçerli promosyon süresi hiçbir zaman kesilmez.

\vspace{0.5em}
\noindent\textbf{Yeni Veritabanı Nesneleri:}

\begin{itemize}[leftmargin=1.5em]
  \item Yeni tablo: \texttt{revenuecat\_webhook\_events}
        (\texttt{revenuecat\_event\_id} UNIQUE; \texttt{processing\_status}:
        received | processed | skipped | failed; redakte payload)
  \item \texttt{listing\_promotion\_purchases} tablosuna yeni sütunlar:
        \texttt{refunded\_at}, \texttt{cancel\_reason}, \texttt{revenuecat\_event\_id}
\end{itemize}

% ── Typecheck ────────────────────────────────────────────────────────────────
\section{Typecheck Durumu}

\begin{tabularx}{\textwidth}{lX}
  \toprule
  \textbf{Paket} & \textbf{Durum} \\
  \midrule
  \texttt{@workspace/api-server} & \textcolor{green!60!black}{\textbf{PASS}} \\
  \texttt{@workspace/mobile} & \textcolor{green!60!black}{\textbf{PASS}} \\
  \texttt{@workspace/mockup-sandbox} & \textcolor{orange}{\textbf{MEVCUT HATA}}
    — \texttt{calendar.tsx}, \texttt{spinner.tsx} (önceden var; bu adımla ilgisiz) \\
  \bottomrule
\end{tabularx}

% ── Aktifleştirme ────────────────────────────────────────────────────────────
\section{Aktifleştirme Gereksinimleri}

RevenueCat webhook'unu canlıya almak için aşağıdaki adımlar gereklidir:

\begin{enumerate}[leftmargin=1.5em]
  \item RevenueCat dashboard $\rightarrow$ \textbf{Integrations $\rightarrow$ Webhooks}
        $\rightarrow$ URL: \texttt{https://\textless domain\textgreater/api/webhooks/revenuecat}
  \item Ortam değişkeni: \texttt{REVENUECAT\_WEBHOOK\_AUTH\_TOKEN}
        (rastgele, güçlü bir string — örn. 32+ karakter)
  \item Opsiyonel (önerilen): \texttt{REVENUECAT\_WEBHOOK\_SIGNING\_SECRET}
        (RevenueCat'in imzalama anahtarı)
\end{enumerate}

% ── Ürün Özellikleri ─────────────────────────────────────────────────────────
\section{Ürün Özellikleri Genel Bakış}

\begin{itemize}[leftmargin=1.5em]
  \item İnteraktif harita — İstanbul'daki sokak hayvanları; renkli markerlar
        (durum: aç / yaralı / sağlıklı / bilinmiyor)
  \item Sokak hayvanı rapor CRUD — fotoğraf, konum, yorum, mama verme
  \item Evcil hayvan profil yönetimi — aşı bilgisi, besleme notları
  \item Sahiplendirme ilan panosu — iletişim bilgisi, başvuru formu
  \item Tam kimlik doğrulama akışı — kayıt / giriş / çıkış / şifre sıfırlama
  \item IAP öne çıkarma sistemi — stacking, ledger tabanlı iade yönetimi
\end{itemize}

\vspace{2em}
\noindent\rule{\textwidth}{0.4pt}
\begin{center}
  \small\color{gray}
  CanYoldaşı — Topluluk Odaklı Sokak Hayvanı Uygulaması \\
  pnpm monorepo · Expo SDK 54 · Express 5 · PostgreSQL · RevenueCat · Stripe
\end{center}

\end{document}
